\documentclass[10pt,a4paper]{article}
\PassOptionsToPackage{dvipsnames}{xcolor}
\usepackage{fullpage}

\usepackage[utf8]{inputenc}

\usepackage[algo2e,ruled,vlined,linesnumbered]{algorithm2e}

\usepackage[normalem]{ulem}
\usepackage[inline]{enumitem}
\usepackage{amsmath,amsfonts,amssymb,amsthm,hyperref,algorithm,array,ragged2e,stmaryrd,bm}
\usepackage{xspace}
\usepackage{tikz}
\usepackage{xcolor}
\usepackage{ifthen}
\usepackage{mathtools} 
\usepackage{bussproofs}
\usepackage{thmtools,thm-restate}

\usepackage{lineno}
\setcounter{secnumdepth}{3} 

\usepackage{thm-restate}
\usepackage{enumitem}
\usepackage[capitalise,nameinlink,noabbrev]{cleveref}

\usepackage[dvipsnames]{xcolor}
\hypersetup{
	colorlinks=true,
	urlcolor=MidnightBlue,
	linkcolor=MidnightBlue,
	citecolor=MidnightBlue,
	unicode
}


% =============== THM Environments =================

%% Theorem environments

\declaretheorem[name=Theorem, parent=section]{theorem}
\declaretheorem[name=Lemma,sibling=theorem]{lemma}
\declaretheorem[name=Claim,sibling=theorem]{claim}
\declaretheorem[name=Proposition,sibling=theorem]{proposition}
\declaretheorem[name=Observation,sibling=theorem]{observation}
\declaretheorem[name=Corollary,sibling=theorem]{corollary}

\declaretheorem[name=Conjecture, sibling=theorem]{conjecture}
\declaretheorem[name=Question, sibling=theorem]{question}
\declaretheorem[name=Open Question, sibling=theorem]{openquest}


\theoremstyle{definition}
\declaretheorem[name=Definition, sibling=theorem, style=definition]{definition}
\declaretheorem[name=Open Problem, sibling=theorem, style=definition]{problem}


\declaretheorem[name=Fact,sibling=theorem]{fact}

\theoremstyle{remark}
\declaretheorem[name=Remark,sibling=theorem,style=remark]{remark}

\makeatletter
\def\th@example{%
  \thm@notefont{}% 
  \normalfont % 
}
\def\th@definition{%
  \thm@notefont{}%
  \normalfont % 
}
\makeatother

\theoremstyle{example}
\newtheorem*{example}{Example}

\usepackage{thmtools,thm-restate}
\theoremstyle{plain}
\declaretheorem[sibling=theorem,name=Theorem]{restatabletheorem}
\declaretheorem[sibling=theorem,name=Lemma]{restatablelemma}

\makeatletter
\renewenvironment{thmt@restatable}[3][]{%
  \thmt@toks{}% 
  \stepcounter{thmt@dummyctr}% 
  \long\def\thmrst@store##1{%
    \@xa\gdef\csname #3\endcsname{%
      \@ifstar{%
        \thmt@thisistheonefalse\csname thmt@stored@#3\endcsname
      }{%
        \thmt@thisistheonetrue\csname thmt@stored@#3\endcsname
      }%
    }%
    \@xa\long\@xa\gdef\csname thmt@stored@#3\@xa\endcsname\@xa{%
      \begingroup
      \ifthmt@thisistheone
        % 
      \else
        \@xa\protected@edef\csname the#2\endcsname{%
          \thmt@trivialref{thmt@@#3}{??} (Restated)}%
        %
        \ifcsname r@thmt@@#3\endcsname\else
          \G@refundefinedtrue
        \fi
        % 
        \@xa\let\csname c@#2\endcsname=\c@thmt@dummyctr
        \@xa\let\csname theH#2\endcsname=\theHthmt@dummyctr
        % 
        \let\label=\@gobble
        \let\ltx@label=\@gobble% 
        % 
        \def\thmt@restorecounters{}%
        \@for\thmt@ctr:=\thmt@innercounters\do{%
          \protected@edef\thmt@restorecounters{%
            \thmt@restorecounters
            \protect\setcounter{\thmt@ctr}{\arabic{\thmt@ctr}}%
          }%
        }%
        % 
        \thmt@trivialref{thmt@@#3@data}{}%
      \fi
      % 
      \ifthmt@restatethis
        \thmt@restatethisfalse
      \else
        \csname #2\@xa\endcsname\ifx\@nx#1\@nx\else[{#1}]\fi
      \fi
      \ifthmt@thisistheone
         % 
         \thmt@rst@storecounters{#3}%
        % 
        \label{thmt@@#3}%
      \fi
      % 
      ##1%
      \csname end#2\endcsname
      % 
      \ifthmt@thisistheone\else\thmt@restorecounters\fi
      \endgroup
    }% 
    \csname #3\@xa\endcsname\ifthmt@thisistheone\else*\fi
    %
    \@xa\end\@xa{\@currenvir}
  }%
  \thmt@collect@body\thmrst@store
}{%
  %% 
}
\makeatother


% =============================================



\newcommand{\deftext}[1]{\emph{#1}}


\renewcommand{\setminus}{\smallsetminus}
\renewcommand{\hat}{\widehat}
\newcommand{\bias}{\mathrm{bias}}
\newcommand{\rk}{\mathrm{rk}}
\newcommand{\s}{\mathrm{s}}
\newcommand{\Ext}{\mathrm{Ext}}
\newcommand{\Sink}{\textsc{Sink}}
\newcommand{\Eq}{\textsc{Eq}}
\newcommand{\Gt}{\textsc{Gt}}
\renewcommand{\And}{\textsc{And}}
\newcommand{\Or}{\textsc{Or}}
\newcommand{\flift}{f\text{-}\mathsf{lift}}
\newcommand{\orlift}{\hspace{-.1em}\lor\hspace{-.1em}\text{-}\mathsf{lift}}
\newcommand{\I}{\operatorname{I}}
\newcommand{\Z}{\mathbb{Z}}
\newcommand{\R}{\mathbb{R}}

\newcommand{\Th}{\mathrm{Th}}
\newcommand{\thtwolift}{\Th(2)\text{-}\mathsf{lift}}
\newcommand{\subcunion}{\textsc{SubcubeUnion}}
\newcommand{\stdist}[2]{\Delta\left(#1,#2\right)}
\newcommand{\entropy}{\mathrm{H}}
\newcommand{\CS}{\textsc{cs}}
\newcommand{\Cell}{\mathsf{C}}
\newcommand{\tribes}{\textsc{Tribes}}
\newcommand{\calS}{\mathcal{S}}
\newcommand{\calR}{\mathcal{R}}
\newcommand{\calG}{\mathcal{G}}

\newcommand{\calC}{\mathcal{C}}
\newcommand{\calU}{\mathcal{U}}
\newcommand{\calD}{\mathcal{D}}
\newcommand{\calT}{\mathcal{T}}
\newcommand{\calF}{\mathcal{F}}
\newcommand{\calQ}{\mathcal{Q}}
\newcommand{\tO}{\tilde{O}}
\newcommand{\NP}{\mathbf{NP}}
\newcommand{\tOmega}{\tilde{\Omega}}
\newcommand{\tTheta}{\tilde{\Theta}}
\DeclareMathOperator{\Exp}{\mathbb{E}}
\DeclareMathOperator{\D}{D}
\newcommand{\cc}{\text{\itshape cc}}
\newcommand{\Dcc}{\D^{\mkern-2mu\cc}}
\newcommand{\Rcc}{\R^{\!\cc}}
\newcommand{\X}{\bm{X}}
\newcommand{\Y}{\bm{Y}}
\newcommand{\K}{\bm{K}}
\newcommand{\x}{\bm{x}}
\newcommand{\y}{\bm{y}}
\renewcommand{\k}{\bm{k}}
\renewcommand{\epsilon}{\varepsilon}
\newcommand{\Good}{\mathcal{T}}
\newcommand{\Xor}{\oplus} % 
\newcommand{\fix}{\textsf{\upshape fix}}
\newcommand{\boundary}{\partial}
\newcommand{\IP}{\mathrm{IP}}
\newcommand{\lb}{\llbracket}
\newcommand{\rb}{\rrbracket}
\newcommand{\bmid}{\,\middle|\,}
\newcommand{\li}{\mathbbm{1}[}
\newcommand{\ri}{]}
\newcommand{\DKL}{\operatorname{D}_{\text{KL}}}
\newcommand{\lpipe}{\bigg|}
\newcommand{\rpipe}{\bigg|}
\newcommand{\poly}{\mathrm{poly}}
\newcommand{\polylog}{\mathrm{polylog}}
\newcommand{\Maj}{\mathrm{Maj}}
\newcommand{\Threshold}{\mathrm{Threshold}}
\newcommand{\Expectation}{\mathbb{E}}
\newcommand{\clique}{\mathrm{CLIQUE}}
\newcommand{\search}{\mathrm{Search}}
\newcommand{\vc}{\mathrm{vc}}
\newcommand{\threshold}{q}

\newcommand{\tribe}{\textsc{Tribes}}
\newcommand{\random}[1]{\bm{#1}}
\newcommand{\newmeasure}[2]{\newcommand{#1}{{\textup{\sffamily #2}}\xspace}}
\newmeasure{\C}{C}

\newcommand{\CoreExtension}{E^*}
\newcommand{\cl}{\mathrm{Cl}}
\newcommand{\CNbound}{R}

\newcommand{\sizeofXi}{\lceil 2^m \ln 2 \rceil}

\newcommand{\ZO}{\{0,1\}}

\newcommand{\BQP}{\textsf{BQP}}
\newcommand{\PH}{\textsf{PH}}
\newcommand{\classP}{\textsf{P}}
\newcommand{\classNP}{\textsf{NP}}

\newcommand{\ACzero}{\textsf{\upshape AC}^0}
\newcommand{\DNF}{\mathrm{DNF}} %
\newcommand{\supp}{\mathrm{supp}}

\newcommand{\Dinv}{\calD^{\text{inv}}}
\newcommand{\im}{\mathrm{Im}}
\newcommand{\cX}{\mathcal{X}}
\newcommand{\cU}{\mathcal{U}}
\newcommand{\cG}{\mathcal{G}}
\newcommand{\Hmin}{\mathrm{H}_{\infty}}
\newcommand{\ind}{\mathbbm{1}}
\newcommand{\xbf}{\bm{x}}
\newcommand{\ybf}{\bm{y}}
\newcommand{\vbf}{\bm{v}}
\newcommand{\Xbf}{\bm{X}}
\newcommand{\Ybf}{\bm{Y}}
\newcommand{\Gbf}{\bm{G}}
\newcommand{\protocolPi}{\Pi}
\newcommand{\protocolPiAlt}{\Pi'}
\newcommand{\protocolPiRand}{\Pi_r}
\newcommand{\Hminof}[1]{\Hmin( {#1} )}
\newcommand{\commCost}[1]{\lvert#1\rvert}
\newcommand{\ta}{t_A}
\newcommand{\tb}{t_B}
\newcommand{\false}{0}
\newcommand{\true}{1}
\newcommand{\treeCP}{\text{CP}^*}
\newcommand{\ResParity}{\mathrm{Res}(\oplus)}

\newcommand{\frege}{\mathrm{Frege^*}}
\newcommand{\res}{\mathrm{Res^*}}


\newcommand{\comment}[1]{}

\ifthenelse{1>0}{
\newcommand{\artur}[1]{\textcolor
{orange}{[Artur: #1]}}
\newcommand{\david}[1]{\textcolor{red}{[David: #1]}}
\newcommand{\duri}[1]{\textcolor{purple}{[Duri: #1]}}
\newcommand{\yassine}[1]{\textcolor{ForestGreen}{[Yassine: #1]}}
\newcommand{\susanna}[1]{\textcolor{blue}{[Susanna: #1]}}
}{
\newcommand{\artur}[1]{}
\newcommand{\david}[1]{}
\newcommand{\duri}[1]{}
\newcommand{\yassine}[1]{}
\newcommand{\susanna}[1]{}
}
\newcommand{\SETSIZE}[1]{\left\lvert#1\right\rvert}
\newcommand{\Setsize}[1]{\bigl\lvert#1\bigr\rvert}
\newcommand{\setsize}[1]{\lvert#1\rvert}
\newcommand{\set}[1]{\{#1\}}


\newcommand{\prflength}{L}
\newcommand{\ResLin}{\textrm{Res}(\oplus)}
\newcommand{\binClique}[2]{\mathsf{BClique}({#1}, {#2})}
\newcommand{\binCliqueGk}{\binClique{G}{k}}
\newcommand{\binCliqueBlock}[2]{\mathsf{BClique}_{\textrm{block}}({#1}, {#2})}
\newcommand{\binCliqueBlockGk}{\binCliqueBlock{G}{k}}

\newcommand{\searchComm}[1]{\mathsf{BClique}^{cc}({#1}, k)}

\newcommand{\binCliqueBlockComm}[2]
{\mathsf{BClique}_{\textrm{block}}^{cc}(#1,#2)}
\newcommand{\binCliqueBlockCommGk}{\binCliqueBlockComm{G}{k}}
\newcommand{\graphSize}{N}

\newcommand{\kClique}{$k$-\textsf{Clique}\xspace}

\providecommand{\erdosrenyi}{Erd\H{o}s-R\'{e}nyi\xspace}
\newcommand{\gnparg}[2]{\mathcal{G}({#1}, {#2})}
\newcommand{\Gnparg}[2]{\mathcal{G}\bigl({#1}, {#2}\bigr)}
\newcommand{\gnp}{\gnparg{n}{p}}
\makeatletter
\newcommand{\aas}{a.a.s.\@\xspace}
\makeatother
\newcommand{\ndense}{almost-complete\xspace}

\newcommand{\F}{\mathbb{F}}

\renewcommand{\algorithmautorefname}{Algorithm}
\renewcommand{\AlgoLineautorefname}{Line}
\renewcommand{\sectionautorefname}{Section}
\renewcommand{\subsectionautorefname}{Section}
\renewcommand{\subsubsectionautorefname}{Section}

\newcommand{\affineDAG}{\Pi}
\newcommand{\edgeProbResLin}{\alpha}
\newcommand{\betabound}{1/300}
\newcommand{\rbound}{1/100(1-\edgeProbResLin)}

\usepackage{fontspec}
\usepackage[margin=22mm]{geometry}
\newboolean{conferenceversion}
\setboolean{conferenceversion}{false}
\newcommand{\myparagraph}[1]{\paragraph{#1}}
